% Part opener. Structural figures get their own number prefix so they
% do not inherit a misleading chapter number.
\structuralfigures{P3.}

\vspace*{0.4cm}
A switch moves frames within a VLAN. Everything that has to cross between VLANs,
between buildings or onto the internet is the router's problem, and this part is
how a router decides.

Domain 3 is 20\% of the exam across four topics, and the verbs are worth noting:
two ask you to \emph{interpret} --- read a routing table, read the status of a
redundancy protocol --- one asks you to \emph{troubleshoot} static routing, and
only one asks you to \emph{configure}, namely OSPF. So this part is as much about
reading output correctly as about typing commands.

\ccnafigh{%
\begin{tikzpicture}[font=\scriptsize]
  \tikzset{
    stg/.style={rounded corners=3pt, draw=#1, fill=#1!8, text=#1!45!black,
                line width=1pt, text width=3.2cm, align=center,
                minimum height=1.55cm, font=\scriptsize\bfseries},
    fl/.style={-{Stealth[length=2.4mm]}, line width=0.9pt, neutral!70}}

  \node[stg=dom3] (a) at (0,0)    {\faIcon{table}\\[2pt]15. The routing table\\
      \footnotesize\mdseries how a router chooses};
  \node[stg=dom3] (b) at (4.1,0)  {\faIcon{map-signs}\\[2pt]16. Static routes\\
      \footnotesize\mdseries you decide the path};
  \node[stg=dom3] (c) at (8.2,0)  {\faIcon{project-diagram}\\[2pt]17--18. OSPF\\
      \footnotesize\mdseries the network decides};
  \node[stg=dom3] (d) at (12.3,0) {\faIcon{shield-alt}\\[2pt]19. HSRP, VRRP\\
      \footnotesize\mdseries the gateway survives};

  \draw[fl] (a) -- (b);
  \draw[fl] (b) -- (c);
  \draw[fl] (c) -- (d);
\end{tikzpicture}}%
{Part~III. Static routing and OSPF are two answers to the same question; first
Chapter~15 establishes how the router picks between whatever answers it has.}{p3map}

\begin{examnote}[One idea underpins the whole part]
A router builds \emph{one} routing table from every source available to it ---
connected interfaces, static routes you typed, and routes learned from OSPF ---
and then forwards on the single best entry for each destination. Administrative
distance decides which \emph{source} wins; the longest prefix match decides which
\emph{entry} wins. Get those two rules straight in Chapter~15 and the rest of the
part is configuration detail.
\end{examnote}
